\chapter{جمع‌بندی، محدودیت‌ها و دفاع}
\label{chap:conclusion}

\section{جمع‌بندی}
این گزارش مسئله مدیریت انرژی ریزشبکه را به‌عنوان یادگیری مقید صورت‌بندی کرد. نشان داده شد که دقت داده‌ای به‌تنهایی کافی نیست، چون تصمیم باید با توازن توان و حدود ذخیره‌ساز سازگار بماند. برای واردکردن این دانش، از چارچوب شبکه عصبی فیزیک‌آگاه استفاده شد \pcite{raissi2019pinns}. برای شناخت متغیرها، تبادل با شبکه و لایه اقتصادی چندذینفع، از صورت‌بندی ریزشبکه متصل به شبکه بهره گرفته شد \pcite{oladejo2019microgrid}.

سهم مشخص پروژه، ساخت یک مدل کاهش‌یافته با ورودی و خروجی صریح، باقیمانده‌های قابل پیاده‌سازی، تابع هزینه چندمؤلفه‌ای و پروتکل ارزیابی است. نظریه بازی به‌عنوان لایه تقسیم سود معرفی شد، نه به‌عنوان نتیجه عددی نسخه‌ فعلی. این مرز، اعتبار علمی گزارش را حفظ می‌کند.

\section{محدودیت‌ها}
\begin{itemize}
  \item مدل الکتریکی مبدل و پایداری گذرا وارد نشده است؛
  \item باتری تک‌حالته و با بازده ثابت است؛
  \item افق تصمیم یک‌گامی است؛
  \item داده واقعی در این نسخه گزارش نشده است؛
  \item لایه چانه‌زنی نش پیاده‌سازی نشده است؛
  \item برتری عددی ادعا نشده است.
\end{itemize}
این محدودیت‌ها نقص پنهان نیستند. هر کدام مسیر کار بعدی را مشخص می‌کنند.

\section{پیشنهاد ادامه}
\begin{enumerate}
  \item ساخت یا گردآوری یک مجموعه‌داده شفاف؛
  \item اجرای جدول~\ref{tab:real-results} با چند بذر؛
  \item بازپارامتری خروجی برای منع شارژ و دشارژ هم‌زمان؛
  \item آموزش روی افق چندگامی؛
  \item افزودن عدم‌قطعیت تابش و قیمت؛
  \item اتصال خروجی انرژی به یک قاعده چانه‌زنی برای چند سایت؛
  \item مقایسه زمان استنتاج با یک حل‌کننده کلاسیک.
\end{enumerate}

\section{چگونه این پروژه را ارائه کنید}
ارائه را روی داستان مسئله بگذارید، نه روی حفظ فرمول‌ها. یک ساختار پانزده تا بیست دقیقه‌ای کافی است:
\begin{enumerate}
  \item یک دقیقه: ریزشبکه چیست و چرا تصمیم آن هم اقتصادی است و هم فیزیکی؛
  \item دو دقیقه: ضعف شبکه فقط‌داده و هزینه حل مکرر بهینه‌سازی؛
  \item سه دقیقه: ایده \lr{PINN} با رابطه~\eqref{eq:pinn-loss} بدون ورود به همه مثال‌های مقاله؛
  \item سه دقیقه: مدل خودتان، شکل~\ref{fig:pinn-workflow} و باقیمانده‌های~\eqref{eq:power-residual} و~\eqref{eq:soc-residual}؛
  \item دو دقیقه: جایگاه مقاله ریزشبکه و اینکه نظریه بازی هنوز لایه بعدی است؛
  \item سه دقیقه: کد، معیارها و جدول خالی به‌عنوان قرارداد صداقت علمی؛
  \item دو دقیقه: محدودیت و کار بعدی.
\end{enumerate}
اسلاید نتیجه جعلی نسازید. اگر آزمایش کوچک دارید، همان را با واحد و بذر نشان دهید. اگر ندارید، پروتکل را نشان دهید و یک اجرای آموزشی روی داده مصنوعی خیلی ساده را زنده اجرا کنید.

\section{پرسش‌های محتمل جلسه دفاع}
\begin{description}
  \item[تفاوت کار شما با مقاله رئیسی چیست؟]
  ایده تابع هزینه فیزیک‌آگاه از آن‌ها است. معادلات، دامنه و هدف این‌جا بهره‌برداری انرژی است، نه حل معادله برگرز یا ناویرـ‌استوکس.
  \item[پس نوآوری شما کجاست؟]
  در انتقال کنترل‌شده روش، تعریف مدل ریزشبکه کاهش‌یافته، جداکردن لایه سود از لایه توان، و گزارش قابل بازتولید.
  \item[چرا عدد ندارید؟]
  چون عدد بدون آزمایش معتبر از عدد نداشتن بدتر است. چارچوب آماده اجرای آن آزمایش است.
  \item[قیدها سخت اعمال می‌شوند؟]
  در نسخه فعلی نرم‌اند. برای قید سخت باید خروجی بازپارامتری شود.
  \item[نظریه بازی را کجا گذاشته‌اید؟]
  بعد از تولید برنامه انرژی، برای تقسیم سود ائتلاف؛ نه داخل گرادیان نسخه فعلی.
  \item[اگر داده کم باشد چه می‌شود؟]
  انتظار نظری این است که مؤلفه فیزیک فضا را محدود کند. این انتظار باید با آزمایش کمبود داده سنجیده شود، نه با حدس در جلسه.
\end{description}

\section{سخن پایانی}
شبکه فیزیک‌آگاه وقتی ارزشمند است که هم معیار یادگیری و هم قانون فیزیکی جدی گرفته شوند. این پروژه آن دو را در یک مسئله واقعی انرژی کنار هم گذاشت و از اغراق عددی پرهیز کرد. نسخه بعدی باید همین چارچوب را با داده و آزمایش پر کند. تا آن روز، دفاع درست یعنی فهم روش، فهم مرز ادعا، و توانایی ادامه کار.
